\exercice{Diagramme des phases (P, T) du naphtalène}{
Le naphtalène (noté \ce{N}) donne lieu aux divers équilibres de changement d'état :

\begin{align*}
\ce{N (l) &<=> N (g)} \quad \ln P_l^\star = 22.76 - \frac{5566}{T} \\
\ce{N (s) &<=> N (l)} \quad \ln P_s^\star = 29.48 - \frac{7935}{T}
\end{align*}

On a donné ci-dessus en \si{\pascal} la pression de la vapeur saturante en présence de liquide ($P_l^\star$) ou de solide ($P_s^\star$) au voisinage du point triple.

\Q Déterminer les coordonnées du point triple.

\Q Dessiner l'allure du diagramme des phases (P, T) à l'aide-de quelques points. Que peut-on dire de la courbe de fusion ? Attribuer les divers domaines.

\Q On réalise une détente isotherme du naphtalène solide à $T = 335 ~ \si{\kelvin}$. Qu'observe-t-on ?

\Q On réalise un échauffement isobare ($P = 2000$ \si{\pascal}) du naphtalène solide à $T=335 ~ \si{\kelvin}$. Qu'observe-t-on ? Représenter la courbe dite d'analyse thermique $T = f(t)$ à
puissance de chauffe constante.}